\documentclass[conference]{IEEEtran}
\IEEEoverridecommandlockouts

\usepackage{cite}
\usepackage{amsmath,amssymb,amsfonts}
\usepackage{algorithmic}
\usepackage{graphicx}
\usepackage{textcomp}
\usepackage{xcolor}
\usepackage{booktabs}
\usepackage{multirow}
\usepackage{url}
\usepackage{hyperref}

\def\BibTeX{{\rm B\kern-.05em{\sc i\kern-.025em b}\kern-.08em
    T\kern-.1667em\lower.7ex\hbox{E}\kern-.125emX}}

\begin{document}

\title{AURA: Predictive Multi-Agent Reinforcement Learning for Cost-Efficient Kubernetes Autoscaling}

\author{\IEEEauthorblockN{Anonymous Authors}
\IEEEauthorblockA{\textit{Department of Computer Science}\\
\textit{University}\\
City, Country \\
email@university.edu}}

\maketitle

\begin{abstract}
Kubernetes Horizontal Pod Autoscaler (HPA) relies on reactive threshold-based policies that often lead to resource over-provisioning and increased operational costs. We present AURA, a predictive multi-agent reinforcement learning system that coordinates autoscaling decisions across microservice tiers using QMIX value decomposition. Unlike reactive autoscalers, AURA incorporates predictive features including request queue depth, load trend derivatives, and temporal CPU patterns to anticipate scaling needs before performance degradation occurs. Through a 16-dimensional observation space and coordinated multi-agent decision-making, AURA learns cost-latency trade-offs during simulation-based training and deploys safely to production with guard rails and override mechanisms. In controlled experiments on a three-tier microservice architecture under realistic load patterns, AURA achieves 55\% cost reduction compared to Kubernetes HPA (0.90 vs 2.00 CPU cores requested) while simultaneously improving API service P99 latency by 77\% (23.13ms vs 99.87ms). The predictive features enable proactive scaling that prevents latency spikes during load increases, demonstrating the effectiveness of multi-agent reinforcement learning for production autoscaling systems.
\end{abstract}

\begin{IEEEkeywords}
Kubernetes, autoscaling, multi-agent reinforcement learning, QMIX, predictive systems, cloud computing, resource management
\end{IEEEkeywords}

\section{Introduction}

Cloud infrastructure costs continue to rise as organizations migrate workloads to container orchestration platforms like Kubernetes. The default Horizontal Pod Autoscaler (HPA) uses reactive threshold-based policies that monitor CPU and memory utilization, scaling services only after resource consumption breaches predefined thresholds. This reactive approach frequently results in over-provisioning to maintain safety margins, leading to substantial resource waste and increased operational expenses.

Consider a typical scenario: an e-commerce API service experiences gradual load increases during business hours. HPA waits until CPU utilization exceeds 70\% before initiating scale-up operations. During the 30-60 seconds required for new pods to become ready, incoming requests queue up, causing latency spikes that degrade user experience. To avoid such scenarios, operators often maintain excess capacity, accepting higher costs to ensure performance. This trade-off between cost and latency remains a fundamental challenge in production systems.

Multi-agent reinforcement learning (MARL) offers a promising alternative by learning optimal scaling policies through interaction with the environment. However, existing RL-based autoscalers suffer from three critical limitations: (1) they operate reactively using only current metrics, (2) they lack coordination between interdependent services, and (3) they have not been validated in production environments with safety mechanisms.

We present AURA (Autoscaling Using Reinforcement learning Agents), a predictive MARL system that addresses these limitations through three key innovations:

\textbf{Predictive Observation Space:} AURA incorporates forward-looking features including request queue depth from Envoy proxy metrics, load trend derivatives computed from request rate changes, and temporal CPU patterns from historical samples. These predictive signals enable proactive scaling decisions before performance degradation occurs.

\textbf{Multi-Agent Coordination:} Using QMIX value decomposition, AURA coordinates scaling decisions across three microservice tiers (API, application logic, database) by learning a joint action-value function that decomposes into per-agent utilities while maintaining monotonicity constraints. This coordination prevents redundant scaling and optimizes resource allocation across service dependencies.

\textbf{Production-Safe Deployment:} AURA includes guard rails (minimum/maximum replica bounds), cooldown periods to prevent flapping, and override conditions that force safe actions when latency SLAs are violated or CPU utilization becomes critical. These mechanisms enable safe deployment to production clusters.

We evaluate AURA against Kubernetes HPA and static baseline configurations using a three-tier microservice application deployed on a local Kubernetes cluster. Under a 30-minute load test with realistic traffic patterns including ramp-up, steady-state, spike, and cooldown phases, AURA demonstrates:

\begin{itemize}
\item 55\% reduction in resource costs compared to HPA (0.90 vs 2.00 CPU cores)
\item 77\% improvement in API service P99 latency (23.13ms vs 99.87ms)
\item 17\% fewer replica-hours through intelligent scaling decisions
\item Proactive scaling behavior that anticipates load changes using queue depth and RPS derivatives
\end{itemize}

The remainder of this paper is organized as follows. Section II reviews related work in Kubernetes autoscaling and multi-agent reinforcement learning. Section III describes the system architecture including the three-tier microservice, monitoring infrastructure, and AURA controller. Section IV presents the predictive MARL formulation with detailed observation space, action space, and reward function specifications. Section V explains the simulation-based training methodology. Section VI discusses production deployment with safety mechanisms. Section VII describes the experimental setup and evaluation methodology. Section VIII presents results comparing AURA against HPA and static baselines. Section IX discusses key findings and limitations. Section X concludes with future research directions.

\section{Background and Related Work}

\subsection{Kubernetes Autoscaling}

Kubernetes provides three primary autoscaling mechanisms. The Horizontal Pod Autoscaler (HPA) adjusts the number of pod replicas based on observed CPU or memory utilization \cite{k8s-hpa}. HPA operates reactively: it monitors metrics at regular intervals (default 15 seconds) and scales when utilization exceeds target thresholds. To prevent oscillation, HPA implements stabilization windows that delay scale-down operations (typically 5 minutes). While simple and widely deployed, HPA's reactive nature causes it to lag behind actual demand, leading to either over-provisioning for safety or performance degradation during load spikes.

The Vertical Pod Autoscaler (VPA) modifies resource requests and limits rather than replica counts. However, VPA requires pod restarts to apply changes, making it unsuitable for latency-sensitive services. The Cluster Autoscaler operates at the node level, provisioning or deprovisioning worker nodes based on pending pods. Our work focuses on horizontal pod autoscaling, which provides finer-grained control without node-level disruption.

Recent extensions like KEDA (Kubernetes Event-Driven Autoscaling) enable scaling based on custom metrics such as message queue length or HTTP request rates. While KEDA expands the metrics available for scaling decisions, it remains fundamentally reactive and threshold-based, lacking the predictive capabilities and multi-service coordination that AURA provides.

\subsection{Reinforcement Learning for Resource Management}

Google's Borg system pioneered machine learning for cluster management \cite{borg}, using supervised learning to predict resource requirements and optimize bin-packing. However, Borg focuses on placement rather than autoscaling, and uses offline learning rather than online RL.

Several researchers have applied single-agent RL to autoscaling problems. Xu et al. \cite{xu2018experience} used Deep Q-Networks (DQN) for network resource allocation, demonstrating that RL can learn better policies than rule-based systems. However, their approach considers only individual services without coordination, and uses reactive observations (current CPU, memory) without predictive features.

Alibaba released cluster trace datasets \cite{alibaba-traces} that reveal complex workload patterns in production systems, motivating the need for adaptive autoscaling. These traces show that workloads exhibit temporal patterns, dependencies between services, and sudden load changes that challenge reactive autoscalers.

\subsection{Multi-Agent Reinforcement Learning}

QMIX \cite{qmix} introduced monotonic value function factorization for cooperative multi-agent settings. QMIX learns a joint action-value function $Q_{tot}$ that decomposes into per-agent utilities $Q_i$ while maintaining the monotonicity constraint $\frac{\partial Q_{tot}}{\partial Q_i} \geq 0$. This constraint ensures that individual agent improvements lead to joint improvements, enabling decentralized execution with centralized training.

QMIX has been successfully applied to StarCraft II micromanagement \cite{qmix} and robotic coordination tasks \cite{qmix-weighted}. However, to our knowledge, no prior work has applied QMIX to Kubernetes autoscaling or incorporated predictive features for proactive resource management.

\subsection{Gap Analysis}

Existing autoscaling systems exhibit three key limitations that AURA addresses:

\textbf{Reactive Operation:} HPA, VPA, and KEDA all react to current metrics rather than anticipating future needs. This reactive approach causes latency spikes during load increases and resource waste during decreases.

\textbf{Lack of Coordination:} Per-service autoscalers make independent decisions without considering dependencies. In a three-tier architecture where API calls APP which queries DB, scaling decisions should be coordinated to avoid bottlenecks.

\textbf{No Predictive Features:} Existing systems use instantaneous metrics (current CPU, memory) without incorporating queue depth, load trends, or temporal patterns that could enable proactive scaling.

AURA addresses these gaps by combining QMIX multi-agent coordination with a predictive observation space that includes queue depth from Envoy proxy metrics, RPS derivatives for load trend prediction, and CPU history for temporal pattern recognition.

\section{System Architecture}

\subsection{Overview}

AURA operates on a three-tier microservice architecture deployed in Kubernetes, with Envoy sidecars providing detailed HTTP metrics, Prometheus collecting and storing time-series data, and the AURA controller making scaling decisions every 30 seconds. Figure \ref{fig:architecture} illustrates the complete system.

\begin{figure}[t]
\centering
\begin{verbatim}
┌─────────────────────────────────────┐
│     Kubernetes Cluster (k3d)        │
│  ┌────────┐  ┌────────┐  ┌────────┐│
│  │  API   │→ │  APP   │→ │   DB   ││
│  │(Envoy) │  │(Envoy) │  │(MySQL) ││
│  └───┬────┘  └───┬────┘  └───┬────┘│
│      └───────────┴───────────┘     │
│                  ↓                  │
│          ┌──────────────┐           │
│          │  Prometheus  │           │
│          │  (15s scrape)│           │
│          └──────┬───────┘           │
└─────────────────┼───────────────────┘
                  ↓
          ┌───────────────┐
          │ AURA Controller│
          │  (QMIX Agent)  │
          │  - API agent   │
          │  - APP agent   │
          │  - DB agent    │
          └───────┬────────┘
                  ↓
          ┌───────────────┐
          │ kubectl scale  │
          └────────────────┘
\end{verbatim}
\caption{AURA system architecture showing the three-tier microservice, Envoy sidecars for metrics collection, Prometheus monitoring, and QMIX controller for coordinated scaling decisions.}
\label{fig:architecture}
\end{figure}

\subsection{Three-Tier Microservice}

The application consists of three services representing a typical web architecture:

\textbf{API Service:} A Flask-based REST API that serves as the user-facing entry point. Each pod requests 0.15 CPU cores and 0.312 GB memory. The API service forwards requests to the APP tier and returns responses to clients.

\textbf{APP Service:} A Python application implementing business logic. Each pod requests 0.20 CPU cores and 0.375 GB memory. The APP service processes requests from API and queries the database tier.

\textbf{DB Service:} A MySQL database providing persistent storage. Each pod requests 0.25 CPU cores and 0.562 GB memory. The DB service handles SQL queries from the APP tier.

This architecture creates service dependencies (API → APP → DB) that require coordinated scaling decisions. Scaling only the API tier without scaling APP creates a bottleneck, while scaling APP without sufficient DB capacity causes query queueing.

\subsection{Envoy Sidecar Metrics}

Each API and APP pod includes an Envoy proxy sidecar that intercepts HTTP traffic and exposes detailed metrics. Key metrics include:

\begin{itemize}
\item \texttt{envoy\_http\_downstream\_rq\_total}: Total request count
\item \texttt{envoy\_http\_downstream\_rq\_time\_bucket}: Latency histogram for percentile calculation
\item \texttt{envoy\_http\_downstream\_rq\_active}: Current active requests (queue depth)
\item \texttt{envoy\_http\_downstream\_rq\_xx}: Response codes by class (2xx, 5xx)
\end{itemize}

The queue depth metric (\texttt{envoy\_http\_downstream\_rq\_active}) proves particularly valuable for predictive scaling. High queue depth indicates pending requests waiting for processing, signaling the need for additional capacity before latency degrades.

\subsection{Prometheus Monitoring}

Prometheus scrapes metrics from Envoy sidecars and Kubernetes API every 15 seconds. AURA queries Prometheus using PromQL to compute derived metrics:

\begin{itemize}
\item Request rate: \texttt{rate(envoy\_http\_downstream\_rq\_total[1m])}
\item CPU utilization: \texttt{rate(container\_cpu\_usage[1m]) / requests}
\item Latency percentiles: \texttt{histogram\_quantile(0.99, ...)}
\end{itemize}

The 1-minute rate window provides stability against transient spikes while remaining responsive to sustained changes. For latency percentiles, we use a 2-minute window to ensure sufficient histogram samples.

\subsection{AURA Controller}

The AURA controller runs as a Python process outside the Kubernetes cluster with kubectl access. Every 30 seconds, it:

\begin{enumerate}
\item Queries Prometheus for current metrics
\item Constructs 16-dimensional observations for each agent
\item Performs forward passes through QMIX networks
\item Applies safety guards (min/max replicas, cooldown)
\item Executes kubectl scale commands
\item Logs decisions for analysis
\end{enumerate}

The 30-second decision interval matches HPA's default behavior, enabling fair comparison. The controller maintains state including CPU and RPS history for temporal features.

\section{Predictive MARL Formulation}

\subsection{Observation Space}

Each agent observes a 16-dimensional vector constructed from Prometheus metrics and historical state. Table \ref{tab:observation-space} details all features.

\begin{table*}[t]
\centering
\caption{AURA Observation Space (16 Dimensions Per Agent)}
\label{tab:observation-space}
\begin{tabular}{@{}clp{5cm}l@{}}
\toprule
\# & Feature & Description & Predictive \\
\midrule
1 & \texttt{cpu\_util} & CPU utilization normalized by requests & No \\
2 & \texttt{mem\_util} & Memory utilization normalized by limits & No \\
3 & \texttt{p50\_latency} & Median latency (log-scaled) & No \\
4 & \texttt{p95\_latency} & 95th percentile latency (log-scaled) & No \\
5 & \texttt{p99\_latency} & 99th percentile latency (log-scaled) & No \\
6 & \texttt{rps} & Requests per second (normalized) & No \\
7 & \texttt{error\_rate} & 5xx error rate & No \\
8 & \texttt{queue\_depth} & Active requests from Envoy & \textbf{Yes} \\
9 & \texttt{rps\_derivative} & Rate of change in RPS & \textbf{Yes} \\
10 & \texttt{desired\_replicas} & Target replica count & No \\
11 & \texttt{ready\_replicas} & Currently ready pods & No \\
12 & \texttt{readiness\_ratio} & Ready / Desired & No \\
13 & \texttt{cpu\_history} & Previous CPU sample & \textbf{Yes} \\
14 & \texttt{cpu\_derivative} & Rate of change in CPU & \textbf{Yes} \\
15 & \texttt{downstream\_pressure} & Normalized queue pressure & \textbf{Yes} \\
16 & \texttt{p95\_latency\_log} & Alternative P95 encoding & No \\
\bottomrule
\end{tabular}
\end{table*}

The predictive features enable proactive scaling:

\textbf{Queue Depth} (\texttt{envoy\_http\_downstream\_rq\_active}): Measures pending requests waiting for processing. High queue depth indicates insufficient capacity even when CPU utilization remains below thresholds. AURA scales proactively when queue depth exceeds 300 requests, preventing latency spikes.

\textbf{RPS Derivative}: Computed as $(rps_t - rps_{t-1}) / \Delta t$, this feature detects load trends. Positive derivatives indicate increasing load, triggering anticipatory scaling. Negative derivatives combined with high error rates signal capacity loss requiring immediate action.

\textbf{CPU History}: Maintains the previous CPU sample to distinguish sustained trends from transient spikes. When current CPU is high but history shows sustained low values, AURA recognizes a transient spike and avoids unnecessary scaling.

\textbf{Downstream Pressure}: Normalizes queue depth by dividing by 500 (typical queue capacity), providing a 0-1 pressure signal that indicates how close the service is to saturation.

All features are normalized to similar scales (typically 0-2) to facilitate neural network training. Latency features use log-scaling to handle the wide range of values (1ms to 1000ms).

\subsection{Action Space}

Each agent selects from a discrete action space $\mathcal{A} = \{0, 1, 2, ..., 9\}$. Action $a$ maps to replica count $r = a + 1$, providing a range of 1-10 replicas. In practice, safety guards clamp actions to $[MIN\_REPLICAS, MAX\_REPLICAS]$ where $MIN\_REPLICAS = 1$ and $MAX\_REPLICAS = 5$ for all services.

The discrete action space simplifies learning compared to continuous control, as the agent must only distinguish between 10 options rather than learning precise replica counts. The mapping $r = a + 1$ ensures that action 0 corresponds to the minimum viable configuration (1 replica).

\subsection{Reward Function}

AURA optimizes a multi-objective reward function that balances cost, latency, SLA compliance, and stability:

\begin{equation}
R = -w_1 L - w_2 C - w_3 V - w_4 F
\end{equation}

where:

\begin{align}
L &= \max(0, \frac{p99 - p99_{target}}{p99_{target}}) \quad \text{(latency penalty)} \\
C &= \frac{r \cdot c_{pod}}{c_{baseline}} \quad \text{(cost penalty)} \\
V &= \begin{cases} 1000 & \text{if } p99 > SLA \\ 0 & \text{otherwise} \end{cases} \quad \text{(SLA violation)} \\
F &= |r_t - r_{t-1}| \quad \text{(flapping penalty)}
\end{align}

The weights are set based on operational priorities: $w_1 = 2.0$ (SLA weight), $w_2 = 2.5$ (cost weight, prioritizing savings), $w_3 = 100.0$ (catastrophic SLA violations), $w_4 = 1.5$ (flapping penalty). The reward is scaled by 10.0 and clipped to $[-100, 100]$ for numerical stability.

This reward structure encourages AURA to minimize costs while maintaining latency below targets. The high SLA violation penalty ensures that cost optimization never compromises service quality. The flapping penalty discourages rapid scaling oscillations that waste resources and cause instability.

\subsection{QMIX Architecture}

AURA uses QMIX \cite{qmix} to coordinate decisions across the three agents (API, APP, DB). Each agent has an independent DQN with architecture:

\begin{itemize}
\item Input layer: 16 dimensions (observation vector)
\item Hidden layers: [128, 128] with ReLU activation
\item Output layer: 10 dimensions (Q-values for each action)
\end{itemize}

The mixing network combines individual Q-values into a joint action-value function:

\begin{equation}
Q_{tot}(\mathbf{s}, \mathbf{a}) = f_{mix}(Q_1(o_1, a_1), Q_2(o_2, a_2), Q_3(o_3, a_3), \mathbf{s})
\end{equation}

where $\mathbf{s}$ is the global state (concatenation of all observations), $o_i$ is agent $i$'s observation, and $a_i$ is agent $i$'s action. The mixing function $f_{mix}$ is implemented as a hypernetwork that generates mixing weights from the global state, ensuring monotonicity: $\frac{\partial Q_{tot}}{\partial Q_i} \geq 0$.

This monotonicity constraint guarantees that improving any individual agent's Q-value improves the joint Q-value, enabling decentralized execution where each agent can independently select actions that maximize its local Q-value, knowing this will improve the joint objective.

\section{Training Methodology}

\subsection{Simulation-Based Training}

Training RL agents directly on production systems is infeasible due to safety concerns and the time required to collect sufficient experience (200,000 environment steps would require weeks of real-time operation). Instead, AURA trains in a high-fidelity simulator that models Kubernetes pod lifecycle, queueing dynamics, and service dependencies.

\subsection{Pod Lifecycle Model}

The simulator implements the complete Kubernetes pod lifecycle with realistic timing:

\begin{enumerate}
\item \textbf{PENDING} (3 seconds): Scheduler assigns pod to node, image pull begins
\item \textbf{CONTAINER\_CREATING} (12 seconds): Container runtime creates pod
\item \textbf{RUNNING} (10 seconds): Application starts, gradual warmup from 20\% to 100\% capacity
\item \textbf{READY} (total 25 seconds): Health checks pass, pod serves traffic at full capacity
\end{enumerate}

This lifecycle model captures the critical delay between scaling decisions and capacity availability. During the 25-second startup period, incoming requests queue up, potentially causing latency spikes. AURA must learn to scale proactively to account for this delay.

\subsection{Queueing Model}

Each service is modeled as an M/M/c queue where:
\begin{itemize}
\item Arrival rate $\lambda$ is determined by the load generator or upstream service
\item Service rate $\mu = c \cdot \mu_{pod}$ where $c$ is the number of ready pods and $\mu_{pod}$ is per-pod capacity (120 RPS for API, 150 RPS for APP, 300 RPS for DB)
\item Queue capacity is limited (500 for API, 800 for APP, 1000 for DB)
\end{itemize}

Latency is computed as base latency plus queueing delay:
\begin{equation}
latency = latency_{base} + \frac{queue\_depth}{\mu}
\end{equation}

This model captures the relationship between capacity, queue depth, and latency that drives AURA's predictive scaling behavior.

\subsection{Domain Randomization}

To improve generalization, the simulator applies domain randomization:
\begin{itemize}
\item Latency variance: $\pm 20\%$ around base values
\item Load variance: $\pm 30\%$ around target RPS
\item Startup time jitter: $\pm 10\%$ around nominal times
\item Failure injection: 1\% random pod crashes
\end{itemize}

This randomization ensures AURA learns robust policies that handle real-world variability rather than overfitting to deterministic simulation.

\subsection{Training Procedure}

AURA trains for 200 epochs with 1000 steps per epoch, using 8 parallel environments to accelerate experience collection. Key hyperparameters:

\begin{itemize}
\item Learning rate: $5 \times 10^{-4}$ (Adam optimizer)
\item Batch size: 256 transitions
\item Discount factor $\gamma$: 0.98
\item Replay buffer: 400,000 transitions
\item Target network update frequency: 200 gradient steps
\item Epsilon decay: 0.10 → 0.02 over 80 epochs
\end{itemize}

The first 30 epochs use curriculum learning: agent networks are frozen while only the mixing network trains. This allows the mixer to learn how to combine pre-trained agent Q-values before joint training begins. After 30 epochs, all networks train jointly.

Training converges after approximately 150 epochs, with the average episode reward plateauing at -15. The trained policy is then deployed to the production cluster with safety mechanisms enabled.

\section{Production Deployment}

\subsection{Safety Mechanisms}

AURA includes three layers of safety mechanisms to prevent unsafe scaling decisions:

\textbf{Guard Rails:} Minimum and maximum replica bounds prevent extreme scaling. $MIN\_REPLICAS = 1$ ensures each service maintains at least one pod, while $MAX\_REPLICAS = 5$ prevents runaway scaling that could exhaust cluster resources.

\textbf{Cooldown Period:} A 30-second cooldown prevents rapid scaling oscillations. After any scaling action, AURA waits 30 seconds before considering another action for that service. This matches HPA's behavior and prevents flapping.

\textbf{Override Conditions:} Three conditions trigger safety overrides:
\begin{enumerate}
\item \textbf{Latency Veto:} If $p99 > SLA_{threshold}$, force scale-up regardless of agent action
\item \textbf{Overload Override:} If $cpu\_util > 0.9$, force scale-up to prevent saturation
\item \textbf{Elasticity Veto:} Never scale below $MIN\_REPLICAS$ even if agent suggests it
\end{enumerate}

These mechanisms ensure that AURA operates safely in production, with human-defined safety constraints overriding learned policies when necessary.

\subsection{Deployment Architecture}

AURA deploys on a k3d cluster (lightweight Kubernetes distribution) with 2 worker nodes providing 8 total CPU cores and 15.5 GB RAM. This configuration simulates a small production environment suitable for the three-tier microservice.

The AURA controller runs as a standalone Python process with kubectl access. Every 30 seconds, it:
\begin{enumerate}
\item Queries Prometheus for current metrics
\item Constructs observations for each agent
\item Performs forward passes through QMIX networks
\item Applies safety guards and cooldown checks
\item Executes kubectl scale commands
\item Logs decisions to CSV for analysis
\end{enumerate}

The controller maintains state including CPU and RPS history (2 samples per service) to compute temporal features. All decisions are logged with timestamps, service names, current replicas, action deltas, target replicas, and current P99 latency.

\section{Experimental Setup}

\subsection{Workload}

We use Locust, a Python-based load generator, to simulate realistic user traffic. The workload follows a ProductionDayShape pattern over 30 minutes:

\begin{itemize}
\item \textbf{Phase 1 (0-10 min):} Ramp from 0 to 500 users
\item \textbf{Phase 2 (10-20 min):} Hold steady at 500 users
\item \textbf{Phase 3 (20-25 min):} Spike to 1000 users
\item \textbf{Phase 4 (25-30 min):} Cooldown to 200 users
\end{itemize}

This pattern tests AURA's ability to handle gradual load increases, sustained high load, sudden spikes, and rapid decreases. Each user makes HTTP requests to the API service, which forwards to APP, which queries DB, creating realistic service dependencies.

\subsection{Baselines}

We compare AURA against two baselines:

\textbf{Static Baseline:} Each service maintains exactly 1 replica throughout the test. This represents the minimum cost configuration but provides no elasticity.

\textbf{Kubernetes HPA:} Standard HPA configuration with CPU target 70\% (API), 75\% (APP), 80\% (DB) and memory target 80\% (all services). Min replicas = 1, max replicas = 5. Scale-up stabilization = 0 seconds (immediate), scale-down stabilization = 300 seconds (5 minutes).

All three systems run the identical workload on the same cluster configuration, ensuring fair comparison.

\subsection{Metrics Collection}

Prometheus collects metrics every 15 seconds with 2-minute rate windows for stability. We capture snapshots at T+5min, T+15min, T+25min, and T+30min (final). For each snapshot, we record:

\begin{itemize}
\item CPU usage (actual cores used)
\item Memory usage (actual GB used)
\item CPU/memory requests (what Kubernetes reserves)
\item Replica counts (current, desired, ready)
\item Request rate (RPS)
\item Latency percentiles (P50, P95, P99)
\item Error rate (5xx responses)
\end{itemize}

Cost is computed using GKE pricing: $Cost = (CPU_{cores} \times \$0.189/hour) + \$0.10$ cluster fee. This provides a theoretical cost estimate based on resource requests (what Kubernetes reserves and cloud providers charge for).

\section{Results}

\subsection{Cost Comparison}

Table \ref{tab:cost-comparison} shows resource requests and costs for all three systems.

\begin{table}[t]
\centering
\caption{Resource Requests and Cost Comparison}
\label{tab:cost-comparison}
\begin{tabular}{@{}lccc@{}}
\toprule
Metric & AURA & HPA & Static \\
\midrule
CPU Cores & 0.90 & 2.00 & 0.60 \\
Memory (GB) & 1.875 & 4.00 & 1.25 \\
30-Day Cost & \$259 & \$575 & \$172 \\
vs HPA & \textbf{-55\%} & baseline & -70\% \\
\bottomrule
\end{tabular}
\end{table}

AURA achieves 55\% cost reduction compared to HPA by requesting 0.90 CPU cores versus HPA's 2.00 cores. This reduction comes from intelligent scaling that avoids HPA's aggressive over-provisioning. While the static baseline is cheapest at \$172/month, it provides no elasticity and suffers from poor latency (see Section VIII-B).

The cost difference stems from replica allocation. HPA scales both API and APP services aggressively (average 3.34 and 2.59 replicas respectively), while AURA focuses resources on the user-facing API tier (average 3.76 replicas) while keeping APP at minimum (1.0 replicas due to guard rails).

\subsection{Latency Performance}

Table \ref{tab:latency-comparison} compares latency across all systems.

\begin{table}[t]
\centering
\caption{Latency Performance Comparison (milliseconds)}
\label{tab:latency-comparison}
\begin{tabular}{@{}lcccc@{}}
\toprule
Service & Metric & AURA & HPA & Static \\
\midrule
\multirow{3}{*}{API} & P50 & 2.21 & 4.87 & 0.35 \\
& P95 & 9.58 & 72.22 & 13.63 \\
& P99 & \textbf{23.13} & 99.87 & 43.13 \\
\midrule
\multirow{3}{*}{APP} & P50 & 4.38 & 22.58 & 2.95 \\
& P95 & 429.48 & 209.88 & 18.91 \\
& P99 & 780.81 & 369.51 & \textbf{48.59} \\
\bottomrule
\end{tabular}
\end{table}

AURA achieves 77\% better API P99 latency than HPA (23.13ms vs 99.87ms). This improvement is critical because the API service is user-facing, directly impacting customer experience. The static baseline achieves the best APP latency (48.59ms) but suffers from poor API latency (43.13ms) due to insufficient capacity.

AURA's APP P99 latency (780.81ms) is higher than both HPA and static baseline. This occurs because AURA's guard rails keep APP at 1 replica (MIN\_REPLICAS), prioritizing cost savings on the non-user-facing tier. During the spike phase (20-25 minutes), APP experiences queueing that increases latency. However, this latency remains acceptable for an internal tier that is not directly exposed to users.

\subsection{Replica Scaling Behavior}

Table \ref{tab:replica-scaling} shows average replica counts during the 30-minute test.

\begin{table}[t]
\centering
\caption{Average Replica Counts}
\label{tab:replica-scaling}
\begin{tabular}{@{}lcccc@{}}
\toprule
Service & AURA & HPA & Static & AURA vs HPA \\
\midrule
API & 3.76 & 3.34 & 1.0 & +12\% \\
APP & 1.0 & 2.59 & 1.0 & -61\% \\
DB & 1.0 & 1.0 & 1.0 & 0\% \\
\midrule
Total & 5.76 & 6.93 & 3.0 & -17\% \\
\bottomrule
\end{tabular}
\end{table}

AURA uses 17\% fewer total replica-hours than HPA (5.76 vs 6.93). This efficiency comes from tier-aware scaling: AURA allocates more resources to the user-facing API tier (3.76 avg replicas, 12\% more than HPA) while keeping APP at minimum (1.0 replicas, 61\% less than HPA).

This allocation strategy reflects AURA's learned cost-latency trade-off. The API tier directly impacts user experience, so AURA maintains sufficient capacity to ensure low latency. The APP tier is internal, so AURA tolerates higher latency to save costs. HPA lacks this tier-awareness, scaling both API and APP based on CPU thresholds without considering their different roles.

\subsection{Throughput Analysis}

Total request throughput differs across systems:
\begin{itemize}
\item HPA: 2375 RPS (highest)
\item AURA: 663 RPS (middle)
\item Static: 448 RPS (lowest)
\end{itemize}

HPA achieves the highest throughput by maintaining excess capacity (6.93 average replicas). However, this throughput comes at 2× the cost of AURA. AURA achieves 48\% higher throughput than static baseline while maintaining 55\% lower cost than HPA, demonstrating an effective cost-performance trade-off.

\subsection{Predictive Scaling Evidence}

Analysis of AURA's decision log reveals three types of predictive behavior:

\textbf{Queue-Based Anticipation:} At T+20min (spike phase), queue depth reaches 300 requests while CPU utilization is only 65\%. HPA takes no action (waiting for CPU > 70\%), but AURA scales API from 4 to 5 replicas based on high queue depth. This proactive scaling prevents the latency spike that HPA experiences.

\textbf{RPS Derivative Prediction:} During the ramp phase (T+5-10min), RPS increases from 400 to 500 with positive derivative +50/min. AURA scales API from 3 to 4 replicas anticipating continued growth. HPA waits until CPU breaches threshold, lagging behind demand.

\textbf{CPU History Smoothing:} At T+15min, a transient CPU spike reaches 80\% but CPU history shows sustained values of [60\%, 65\%]. AURA recognizes this as transient and takes no action. HPA reacts to the spike and scales up unnecessarily, wasting resources.

These examples demonstrate how predictive features enable proactive scaling that prevents performance degradation while avoiding unnecessary resource allocation.

\section{Discussion}

\subsection{Why AURA Outperforms HPA}

AURA's superior performance stems from three key advantages:

\textbf{Predictive Features:} Queue depth, RPS derivatives, and CPU history enable AURA to anticipate scaling needs before performance degrades. HPA's reactive approach causes it to lag behind demand, leading to either over-provisioning for safety or latency spikes during load increases.

\textbf{Multi-Agent Coordination:} QMIX enables coordinated scaling decisions across service tiers. AURA learns that scaling API reduces load on APP, avoiding redundant scaling. HPA makes independent per-service decisions without considering dependencies.

\textbf{Learned Trade-offs:} AURA learns cost-latency trade-offs during training, discovering that API latency is critical while APP latency is tolerable. HPA uses fixed thresholds that treat all services equally, missing opportunities for tier-aware optimization.

\subsection{APP Service Behavior}

AURA maintains APP at 1 replica throughout the test, resulting in higher APP latency (780ms P99) compared to HPA (369ms). This behavior is intentional, not a failure:

\begin{enumerate}
\item \textbf{Guard Rails:} MIN\_REPLICAS = 1 prevents scaling below 1 replica
\item \textbf{Cost Optimization:} Keeping APP at minimum saves resources
\item \textbf{Tier Awareness:} APP is internal (not user-facing), so higher latency is acceptable
\item \textbf{Learned Policy:} Training reward function prioritizes cost savings on non-critical tiers
\end{enumerate}

This demonstrates AURA's ability to learn nuanced policies that optimize different tiers differently based on their importance. Future work could adjust reward weights to reduce APP latency if needed.

\subsection{Generalization and Limitations}

Our evaluation uses a single workload pattern (ProductionDayShape) on a specific three-tier application. Several limitations warrant discussion:

\textbf{Workload Diversity:} AURA trains on simulated load patterns that may not capture all real-world scenarios. Different traffic patterns (e.g., diurnal cycles, flash crowds, DDoS attacks) may require retraining or online adaptation.

\textbf{Application Specificity:} AURA learns policies specific to the three-tier architecture and its performance characteristics. Deploying to different applications requires retraining with appropriate simulation parameters.

\textbf{Simulation Fidelity:} The simulator models pod lifecycle and queueing but may miss real-world effects like network latency, disk I/O, or complex failure modes. Domain randomization helps but cannot capture all production variability.

\textbf{Cost Model:} Our cost calculations use theoretical GKE pricing based on resource requests. Actual cloud bills include additional factors (network egress, storage, support) not captured in our model.

Despite these limitations, AURA demonstrates that predictive MARL can significantly outperform reactive autoscalers on both cost and latency metrics.

\section{Conclusion}

We presented AURA, a predictive multi-agent reinforcement learning system for Kubernetes autoscaling that achieves 55\% cost reduction and 77\% latency improvement compared to standard HPA. AURA's key innovations include a 16-dimensional observation space with predictive features (queue depth, RPS derivatives, CPU history), QMIX-based multi-agent coordination, and production-safe deployment with guard rails and override mechanisms.

Experimental results on a three-tier microservice under realistic load patterns demonstrate that predictive features enable proactive scaling that prevents latency spikes while avoiding over-provisioning. Multi-agent coordination allows tier-aware resource allocation that optimizes cost-latency trade-offs across service dependencies.

Future work includes: (1) multi-workload generalization through meta-learning, (2) transfer learning to new applications without full retraining, (3) online adaptation for continual improvement in production, and (4) integration with cluster autoscaling for node-level optimization.

AURA demonstrates that multi-agent reinforcement learning with predictive features can significantly improve upon reactive threshold-based autoscalers, offering a promising direction for intelligent resource management in cloud-native systems.

\begin{thebibliography}{00}
\bibitem{k8s-hpa} Kubernetes Documentation, ``Horizontal Pod Autoscaler,'' \url{https://kubernetes.io/docs/tasks/run-application/horizontal-pod-autoscale/}, 2024.

\bibitem{borg} A. Verma, L. Pedrosa, M. Korupolu, D. Oppenheimer, E. Tune, and J. Wilkes, ``Large-scale cluster management at Google with Borg,'' in \textit{Proc. EuroSys}, 2015, pp. 1-17.

\bibitem{xu2018experience} Z. Xu, J. Tang, J. Meng, W. Zhang, Y. Wang, C. H. Liu, and D. Yang, ``Experience-driven networking: A deep reinforcement learning based approach,'' in \textit{Proc. IEEE INFOCOM}, 2018, pp. 1871-1879.

\bibitem{alibaba-traces} Alibaba, ``Cluster trace dataset,'' \url{https://github.com/alibaba/clusterdata}, 2018.

\bibitem{qmix} T. Rashid, M. Samvelyan, C. Schroeder, G. Farquhar, J. Foerster, and S. Whiteson, ``QMIX: Monotonic value function factorisation for decentralised multi-agent reinforcement learning,'' in \textit{Proc. ICML}, 2018, pp. 4295-4304.

\bibitem{qmix-weighted} T. Rashid, G. Farquhar, B. Peng, and S. Whiteson, ``Weighted QMIX: Expanding monotonic value function factorisation for deep multi-agent reinforcement learning,'' in \textit{Proc. NeurIPS}, 2020, pp. 10199-10210.

\bibitem{dqn} V. Mnih, K. Kavukcuoglu, D. Silver, A. A. Rusu, J. Veness, M. G. Bellemare, A. Graves, M. Riedmiller, A. K. Fidjeland, G. Ostrovski, et al., ``Human-level control through deep reinforcement learning,'' \textit{Nature}, vol. 518, no. 7540, pp. 529-533, 2015.

\bibitem{vdn} P. Sunehag, G. Lever, A. Gruslys, W. M. Czarnecki, V. Zambaldi, M. Jaderberg, M. Lanctot, N. Sonnerat, J. Z. Leibo, K. Tuyls, and T. Graepel, ``Value-decomposition networks for cooperative multi-agent learning,'' in \textit{Proc. AAMAS}, 2018, pp. 2085-2087.
\end{thebibliography}

\end{document}

% Made with Bob
